\section{\filetotitle}\label{sec:variational_principles_eom}
Having introduced the variational principles in the previous section, we now apply them to derive equations of motion for specific variational ansatz. We will focus on various Gaussian-based ansatz, which are widely used in quantum dynamics simulations due to their computational efficiency and ability to capture essential features of wavepacket evolution. We will also discuss the extension of these methods to non-adiabatic dynamics, where multiple electronic states are involved.
\subsection{Single Gaussian Wavepacket Ansatz}
We start with the simplest case of a single Gaussian wavepacket, which we will denote as $\ket{\Psi(t)}=\ket{G(\symbf{\alpha}(t))}$, where $\symbf{\alpha}(t)$ represents the set of parameters defining the Gaussian (e.g., center position, momentum, width). The gaussian wavepacket we will use can be expressed as
\begin{equation}
\ket{G(\symbf{\alpha}(t))}=N\exp\left[\sum_j^D\left(-\frac{1}{2}\gamma_j(t)\left(R_j-q_j(t)\right)^2+\i p_j(t)\left(R_j-q_j(t)\right)\right)+\i\phi(t)\right],
\end{equation}
where $N$ is the normalization factor, $\symbf{\alpha}(t)$ include the center position $\symbf{q}(t)$, momentum $\symbf{p}(t)$, the complex width $\symbf{\gamma}(t)$ and the global phase factor $\phi(t)$. Should we consider the frozen Gaussian approximation, the parameter $\symbf{\gamma}(t)$ is kept constant in time, while in the thawed Gaussian approximation, it is allowed to evolve. By applying the \acrshort{dfvp} to this ansatz, we can derive equations of motion for the parameters $\symbf{\alpha}(t)$, which govern the evolution of the Gaussian wavepacket.
\subsubsection{Frozen Gaussian Evolution}
The first derived equations of motion will be for the frozen Gaussian, which means that the width of the gaussian wavepacket is fixed ($\dot{\symbf{\gamma}}(t)=0$). We will use the \acrshort{dfvp} in Eq.~\ref{eq:dfvp_orthogonality_condition} to derive the equations of motion for the remaining parameters $\symbf{q}(t)$, $\symbf{p}(t)$ and $\phi(t)$. We will also restrict ourselves to the case of a single nuclear degree of freedom for simplicity, but the generalization to multiple dimensions is straightforward.

To apply the \acrshort{dfvp}, we first need to compute all the necessary derivatives of the Gaussian wavepacket with respect to the parameters as
\begin{align}
\ket{v_\phi}&:=\frac{\partial}{\partial \phi}\ket{\Psi}=\i\ket{\Psi},\\
\ket{v_q}&:=\frac{\partial}{\partial q}\ket{\Psi}=\left(\gamma\left(R-q\right)-\i p\right)\ket{\Psi},\\
\ket{v_p}&:=\frac{\partial}{\partial p}\ket{\Psi}=\i\left(R-q\right)\ket{\Psi}.
\end{align}
These vectors represent the tangent space of the variational manifold defined by the Gaussian wavepacket. The \acrshort{dfvp} states that the projection of the time derivative of the wavefunction onto this tangent space must vanish. Let's start with the projection onto $\ket{v_\phi}$ as
\begin{align}
\braket{v_\phi|\dot\Psi+\i H\Psi}&=0\\
-\i\braket{\Psi|\dot\Psi}+\braket{\Psi|H|\Psi}&=0\\
-\i\braket{\gamma(R-q)\dot q+\i\dot p(R-q)-\i p\dot q+\i\dot\phi}&=-\braket{H}\\
\dot\phi&=p\dot q-\braket{H},\label{eq:gaussian_eom_phi}
\end{align}
where we have used the normalization condition $\braket{\Psi|\Psi}=1$ and the fact that $\braket{R-q}=0$. This gave us the equation of motion for the global phase $\phi(t)$, which corresponds to the classical action along the trajectory defined by $q(t)$ and $p(t)$. Next, we project onto $\ket{v_p}$ as
\begin{align}
\braket{v_p|\dot\Psi+\i H\Psi}&=0\\
\braket{\i(R-q)\Psi|\dot\Psi}&=-\braket{(R-q)\Psi|H|\Psi}\\
-\i\braket{\gamma(R-q)^2\dot q+\i\dot p(R-q)^2}&=-\braket{(R-q)H}\\
(\dot p-\i\gamma\dot q)\braket{(R-q)^2}&=-\braket{(R-q)H}.\label{eq:gaussian_eom_p}
\end{align}
Since we are working with Gaussian wavefunction, we can evaluate the expectation values analytically. In particular, we have $\braket{(R-q)^2}=\frac{1}{2\Re\gamma}$. Additionally, we can express the right-hand side, using the full kinetic energy operator, as
\begin{equation}
\braket{(R-q)H}=-\frac{1}{2m}\braket{(R-q)\partial_R^2}+\Braket{(R-q)V}=\frac{\i\gamma p}{m}\braket{(R-q)^2}+\braket{(R-q)V}.
\end{equation}
Plugging this back into Eq.~\ref{eq:gaussian_eom_p}, we get
\begin{equation}
\dot p-\i\gamma\dot q=-\frac{\i\gamma p}{m}-\frac{\braket{(R-q)V}}{\braket{(R-q)^2}}=-\frac{\i\gamma p}{m}-\Braket{\frac{\partial V}{\partial R}}.
\end{equation}
where we have uset the Stein's lemma to evaluate $\braket{(R-q)V}$. Now, we would like to separate the equation into one equation for $\dot p$ and one for $\dot q$. Since $q$ and $p$ are real-valued parameters, we can take the real and imaginary parts of the above equation to get
\begin{align}
\dot q&=\frac{p}{m},\\
\dot p&=-\Braket{\frac{\partial V}{\partial R}}.
\end{align}
These final equations of motion, along with Eq~\eqref{eq:gaussian_eom_phi}, correspond to the classical Hamilton's equations for a particle moving in the potential $V(R)$, where the force is given by the expectation value of the gradient of the potential. This is a remarkable result, as it shows that the evolution of a frozen Gaussian wavepacket follows classical mechanics, with quantum effects encoded in the width and phase of the wavepacket. One can also project onto $\ket{v_q}$ to derive an equation for $\dot q$, but it will yield the same result as above, confirming the consistency of the equations of motion.
\subsubsection{Thawed Gaussian Evolution}
To improve the accuracy of the Gaussian wavepacket, we can allow the width parameter $\gamma(t)$ to evolve in time. This leads to the thawed Gaussian approximation, where the wavepacket can change its shape as it propagates. The equations of motion for the thawed Gaussian can be derived in a similar manner as for the frozen Gaussian, but now we also need to consider the projection onto $\ket{v_\gamma}$, which is given by
\begin{equation}\label{eq:gaussian_derivative_gamma}
\ket{v_\gamma}:=\frac{\partial}{\partial \gamma}\ket{\Psi}=-\frac{1}{2}(R-q)^2\ket{\Psi}.
\end{equation}
We now require that the projection onto $\ket{v_\gamma}$ also vanishes, along with the projections onto $\ket{v_\phi}$, $\ket{v_q}$ and $\ket{v_p}$. Projecting onto $\ket{v_q}$ and $\ket{v_p}$ will yield the same equations of motion for $q(t)$ and $p(t)$ as in the frozen case, since the width parameter introduces only $\Braket{(R-q)^3}$ terms, which are zero. However, the projection onto $\ket{v_\phi}$ will now include an additional term and will yield
\begin{align}
\braket{v_\phi|\dot\Psi+\i H\Psi}&=0\\
\braket{\Psi|\dot\Psi}&=-\i\braket{H}\\
\Braket{-\frac{1}{2}\dot\gamma(R-q)^2+\gamma(R-q)\dot q+\i\dot p(R-q)-\i p\dot q+\i\dot\phi}&=-\i\braket{H}\\
-\frac{1}{2}\dot\gamma\braket{(R-q)^2}-\i p\dot q+\i\dot\phi&=-\i\braket{H}\\
\i\dot\phi&=\i p\dot q+\frac{1}{2}\dot\gamma\braket{(R-q)^2}-\i\braket{H}\\
\dot\phi&=p\dot q-\frac{\i}{4\Re\gamma}\dot\gamma-\braket{H}.
\end{align}
As we can see, the equation of motion for the global phase $\phi(t)$ now includes an additional term that depends on the time derivative of the width parameter $\gamma(t)$. This term accounts for the change in the shape of the wavepacket as it evolves. Finally, let's derive the equation of motion for $\gamma(t)$ by projecting onto $\ket{v_\gamma}$ as
\begin{align}
\braket{v_\gamma|\dot\Psi+\i H\Psi}&=0\\
\Braket{(R-q)^2\Psi|\dot\Psi}&=-\i\Braket{(R-q)^2H}\\
\Braket{-\frac{1}{2}\dot\gamma(R-q)^4+(R-q)^2(\i\dot\phi-\i p\dot q)}&=-\i\Braket{(R-q)^2H}\\
\dot\gamma\Braket{(R-q)^2}^2&=\i\left(\Braket{(R-q)^2H}-\Braket{(R-q)^2}\Braket{H}\right)
\end{align}
where we have used the fact that $\Braket{(R-q)^4}=3\Braket{(R-q)^2}^2$ for a Gaussian wavefunction. The final step is now calculating the covariance term $\Braket{(R-q)^2H}-\Braket{(R-q)^2}\Braket{H}$. Splitting the Hamiltonian into kinetic and potential energy, we can evaluate the kinetic part as
\begin{align}
\braket{T}&=-\frac{1}{2m}\Braket{\partial_R^2}=\frac{1}{2m}\left(\gamma-\gamma^2\Braket{(R-q)^2}+p^2\right)\\
\Braket{(R-q)^2T}=-\frac{1}{2m}\Braket{(R-q)^2\partial_R^2}&=\frac{1}{2m}\left(\gamma\Braket{(R-q)^2}-\gamma^2\Braket{(R-q)^4}+p^2\Braket{(R-q)^2}\right)\\
\Braket{(R-q)^2T}-\Braket{(R-q)^2}\Braket{T}&=-\frac{\gamma^2}{2m}\left(\Braket{(R-q)^4}-\Braket{(R-q)^2}^2\right)=-\frac{\gamma^2}{m}\Braket{(R-q)^2}^2.
\end{align}
As for the potential energy part, we can express it as
\begin{align}
\Braket{(R-q)^2V}-\Braket{(R-q)^2}\Braket{V}&=\Braket{(R-q)^2}^2\Braket{\frac{\partial^2 V}{\partial R^2}},
\end{align}
which can be easily derived using the Stein's lemma. Putting everything together, we finally get the equation of motion for the width parameter $\gamma(t)$ as
\begin{equation}
\dot\gamma=-\frac{\i\gamma^2}{m}+\i\Braket{\frac{\partial^2 V}{\partial R^2}}
\end{equation}
which shows that the width of the Gaussian wavepacket evolves according to a balance between the kinetic energy, which tends to spread the wavepacket, and the curvature of the potential energy surface, which can either focus or defocus the wavepacket depending on its sign.
\subsection{Variational Multiconfigurational Gaussian (vMCG) Approach}
\subsection{Extension to Non-Adiabatic Dynamics: The Single-Set Formulation}
